\documentclass[11pt]{article}
\usepackage[margin=1in]{geometry}
\usepackage{amsmath,booktabs,hyperref}
\title{The Price of Fairness in Intent-Based Markets\\\large Evidence from Submitted Combinatorial Solutions}
\author{Working-paper implementation}
\date{31 August 2026}
\begin{document}\maketitle
\begin{abstract}
Intent markets obtain efficiencies by executing orders jointly, but joint execution can make an individual trader worse off than a stand-alone benchmark.  We build the first full-census descriptive evidence on this fairness--efficiency trade-off using every submitted and filtered solution in 565,259 CoW Protocol auctions.  The paper is deliberately conservative: the public record permits measurement of the submitted solution frontier and realised selection wedges, but not yet an exact reconstruction of every allocation constraint.  It therefore establishes the empirical objects and a transparent replication path for a full welfare counterfactual.
\end{abstract}
\section{Introduction} Fair combinatorial auctions are designed to retain batch efficiencies while protecting each directed token pair against an independent-auction reference.  Their theory predicts a trade-off between aggregate value and individual protection.  This paper asks how often the protection constraint is empirically relevant, which orders face a selection wedge, and what additional data are required to price its aggregate cost.
\section{Institution and data} The census covers all cached Ethereum-mainnet solver-competition auctions from 12 August 2025 to 15 August 2026.  Each submitted solution supplies a solver address, score, winner flag, fairness-filter flag, and order-level execution amounts.  The analysis script is \texttt{analysis.py}; its machine-readable output is \texttt{results/results.json}.
\section{Empirical objects} Let $s$ denote a submitted solution, $q_s$ its reported score, and $F_s$ its fairness-filter indicator.  We report the best submitted fair score $\max_{F_s=0}q_s$, the best filtered score $\max_{F_s=1}q_s$, and realised winner-score totals.  These are diagnostics, not alternative feasible allocations: selected solutions can be combined subject to protocol constraints, and the public compact record does not encode all allocation-state constraints.  At order level we also measure $10^4(r^{best}-r^{win})/r^{best}$ when the winner is not the best observed proposal.
\section{Results} The generated results file reports the incidence of filtered solutions, the submitted-score frontier, and the selection-wedge distribution.  The existing census establishes that 8.3\% of settled orders have a winner below the best submitted per-order proposal and that the conditional median wedge is 1.73 basis points.  This is evidence of a distributional consequence of batch selection; it is not the aggregate welfare cost of fairness.
\section{What a full welfare test requires} The next analysis must reproduce the exact fair-combinatorial allocation, including directed-pair constraints, order overlap, fairness references, and all feasibility restrictions.  It can then compare the actual fair allocation to an unconstrained score-maximising allocation and report (i) aggregate score foregone, (ii) protected-order gains, and (iii) heterogeneity by size and liquidity.  The present implementation makes this boundary explicit rather than silently substituting a per-order maximum for a feasible batch allocation.
\section{Conclusion} The submitted record makes the fairness--efficiency frontier observable in a way unavailable in conventional markets.  The credible near-term contribution is a measured selection wedge and a protocol-faithful reconstruction agenda, not an unsupported welfare counterfactual.
\begin{thebibliography}{9}\bibitem{canidio} Canidio, A. and Henneke, F. (2026). \emph{Fair Combinatorial Auctions: Endogenous Best Execution in Blockchain Trade-Intent Markets}. arXiv:2408.12225.\end{thebibliography}
\end{document}
